\documentclass{ximera}
% xmPreamble.tex en xmPrintstyle.sty worden automatisch geladen door ximera.cls (v2.7.8);
% NIET nog eens \input{./preamble.tex} of \addPrintStyle{.} doen (geeft 'already defined'-fouten).

\begin{document}
\author{Alexander Holvoet}
\xmtitle{Zandhopen: inleiding}{}

Deze reeks lesactiviteiten hoort bij het Uitwiskeling-artikel
\href{https://www.uitwiskeling.be/2026/04/16/zandhopen/}{``Moet er nog zand zijn?''}
(Uitwiskeling 42/2). In dat artikel lees je het volledige didactische verhaal achter de
activiteiten: het ontwerpproces, de leerkrachtenadviezen en verdere inspiratie.

\subsection*{Waar gaat het over?}

Door zand (of zout, rijst, \ldots) te gieten op verschillende basisvormen ontdek je hoe
meetkunde opduikt in een heel concreet experiment. Telkens werk je op dezelfde
onderzoekende manier: je \emph{voorspelt} eerst wat je verwacht, je \emph{voert} dan het
experiment uit, en je \emph{verklaart} ten slotte wat je ziet.

\subsection*{Leerdoelen}

Doorheen de activiteiten ontdek en gebruik je onder andere:
\begin{itemize}
    \item de \emph{rusthoek} van een hoop materiaal en hoe je die meet;
    \item \emph{hoogtelijnen} en de \emph{kamlijn} als meetkundige plaatsen van punten op
        gelijke afstand van de randen;
    \item \emph{bissectrices} en de ingeschreven cirkel van een veelhoek, en de
        bissectricestelling;
    \item \emph{kegelsneden} \textemdash{} cirkel, ellips en parabool \textemdash{} die als
        kamlijn verschijnen, samen met hun definitie als meetkundige plaats en de
        weerkaatsingseigenschap van de ellips;
    \item het verband tussen een ruimtelijke situatie (de zandhoop) en haar bovenaanzicht.
\end{itemize}

\subsection*{Voor wie?}

De activiteiten zijn bedoeld voor leerlingen van het secundair onderwijs. De eerste
paragrafen (rusthoek, hoogtelijnen, bissectrices) zijn breed toegankelijk; de paragrafen
over kegelsneden sluiten aan bij de leerstof van de derde graad. Je kunt de reeks dus
gedeeltelijk of volledig gebruiken, afhankelijk van het niveau.

\subsection*{Wat heb je nodig?}

Droog zand, zout of rijst, en enkele basisvormen uit karton (rond, vierkant, rechthoekig,
driehoekig, \ldots). Plaats je basisvorm wat hoger (bijvoorbeeld op een omgekeerd bekertje),
zodat het teveel aan zand er langs de zijkant kan afglijden. Voor enkele opgaven is
GeoGebra handig om te meten en te construeren.

\end{document}
